\documentclass[11pt]{article}
\usepackage{amsmath,amssymb}
\usepackage{booktabs}
\usepackage[margin=2.3cm]{geometry}

\title{Formato total Python $\to$ \LaTeX{} con \texttt{tex()}}
\author{}
\date{\today}

\begin{document}
\maketitle

La función \verb|tex(obj, prec, env)| convierte cualquier objeto de Python en
\LaTeX{} listo para el documento; \verb|texesc(s)| escapa texto plano de forma
segura. Así el formato lo decides tú desde el propio documento.

%#python
import numpy as np
import pandas as pd
import sympy as sp

A = np.array([[2.0, 1.0], [3.0, 4.0]])
autovalores = np.linalg.eigvals(A)

x = sp.symbols("x")
integral = sp.integrate(sp.sin(x)**2, x)

tabla = pd.DataFrame({
    "Pieza": ["Viga", "Pilar", "Zapata"],
    "L (m)": [6.0, 3.2, 1.5],
    "Coste": [1250.5, 890.25, 410.0],
})

v = ureg.Quantity(25.456, "m/s")
nota = "Coste del 50% & margen #1"
%#end

\section{Matrices (numpy)}
\[ A = \py{tex(A, prec=3)} , \qquad
   \lambda = \py{tex(autovalores, prec=4, env="pmatrix")} \]

\section{Cálculo simbólico (sympy)}
\[ \int \sin^2 x \, dx = \py{tex(integral)} + C \]

\section{Tabla (pandas)}
\begin{center}
\py{tex(tabla, prec=4)}
\end{center}

\section{Unidades (pint) y texto seguro}
La velocidad es $v = \py{tex(v, prec=2)}$ y la nota dice:
«\py{texesc(nota)}».

También puedes formatear a mano con f-strings: \py{f"{float(v.magnitude):.1f}"} m/s.

\end{document}
